\section{Висновки}\label{sec:vysnovky}

\subsection{Підсумок вимог та критерії готовності до реалізації}\label{sec:vysnovky-pidsumok}
Ця специфікація вимог до програмного забезпечення (SRS) фіксує узгоджений консенсусний простір знань $K^*$ для Cinema Booking System, сформований у ході ітеративного сократичного діалогу між автором проєкту та автономним AI-агентом Google Antigravity~\sked{K6,K10}.

Специфікація повністю покриває всі критичні аспекти функціонування системи:
\begin{itemize}
    \item Чітке розмежування ролей Користувача та єдиного системного Адміністратора;
    \item Загальнодоступний перегляд сеансів та захищене створення/скасування бронювань виключно для зареєстрованих клієнтів;
    \item Строге дотримання обмеження від 1 до 6 місць на одне бронювання;
    \item Надійну атомарну модель захисту від подвійних бронювань без застосування складних таймаутів (hold) у версії 1;
    \item Правило «все або нічого» при бронюванні групи місць;
    \item Однозначні правила скасування бронювань користувачем (лише повне скасування суто до часу початку сеансу);
    \item Правила модифікації розкладу адміністратором (блокування редагування сеансів з активними бронюваннями та можливість явного скасування із зазначенням обов'язкової причини);
    \item Обов'язковість дисципліни Test-Driven Development (TDD) як засадничого процесу розробки нової функціональності та усунення дефектів~\sked{K7};
    \item Регламент операційного та тестового логування із гарантією захисту персональних і конфіденційних даних.
\end{itemize}

Зафіксований набір вимог та BDD-сценаріїв у Додатку~\ref{sec:dodatok-matrytsia} є вичерпним і достатнім для переходу до етапу інженерного проєктування та безпосередньої реалізації коду системи за методологією TDD~\sked{K7,K10}.
